\addcontentsline{toc}{section}{Заключение}

В ходе проведённого исследования был осуществлён сравнительный анализ современных технологий разработки интеллектуальных систем: языков представления информации и форматов представления знаний.

В первом разделе работы были подробно рассмотрены языки представления информации Prolog и LINQ. Проанализированы их синтаксические и семантические особенности, выявлены общие черты и различия в подходах к представлению и обработке знаний. Prolog представляет собой классический язык логического программирования, основанный на логике предикатов, в то время как LINQ предоставляет технологию интегрированных запросов в рамках платформы .NET. Результаты анализа были формализованы с помощью SCn-кода.

Во втором разделе исследованы современные инструменты и платформы для работы с онтологиями и знаниями на примере Protégé и Google Knowledge Graph. Protégé представляет собой открытый инструмент для создания и редактирования онтологий, разработанный Стэнфордским университетом. Google Knowledge Graph --- проприетарная платформа представления знаний в виде графа сущностей, используемая в поисковой системе Google. Проведен анализ их функциональных возможностей, методов представления знаний, а также особенностей архитектуры. Информация о каждом инструменте также была формализована на SCn-коде.

Результаты работы могут быть использованы при проектировании гибридных интеллектуальных систем, сочетающих возможности логического вывода, работы с данными и онтологического моделирования, а также при выборе и обосновании технологических решений для конкретных прикладных задач.
